%\input ../util/bookmacros
%\input ../util/docparams

\appendix{Appendix C}{Infinite Sequences and Limits}
Every derivative in this book was computed by the {\it approximate,
then take a limit\/} move, and each time the ``limit'' step was taken
on faith. In this appendix we pay that debt by pinning down what it
means for a sequence of approximations to head toward a value.
That is, we will invent a definition for what we've been calling
the \index{limit} of a \index{sequence}. The definition will come from looking at examples.
I hope these examples will convince you that the definition given 
below is reasonable. 

First to make life easier when discussing sequences let us define the 
{\it tail\/}, $T_N$, of a sequence $\{x_n\}_{n=1}^{\infty}$ at $N$ 
to be the set $\{\, x_n \mid n > N\, \}$.
You can think about forming the tail by visualizing the sequence as a list of
numbers written out sequentially from left to right and then
chopping the sequence at a location $N$; this leaves the 
list $[x_1 \ldots x_N]$ and the rest which we are calling the tail at $N$, 
$T_N$. 

\example{Let $x_n = 1/n$. The tail of this sequence at $10$, $T_{10}$, is 
         $\{\, x_n \mid n > 10\,\} = \{1/11,\; 1/12,\; 1/13, \; \ldots\}$.}

You should draw a picture of the tail.

Notice that the sequence in the example seems to be ``headed toward'' $0$. 
Question: Do you believe this, if so, how do we say this precisely?

We need to gain some experience with sequences in order to 
extract the essence of the notion of ``heading towards''.
To this end, 
let's look at some examples to help us sharpen our idea of what 
this ``heading toward'' or ``goes to''
might mean. We will use the term {\it convergence\/} to 
mean that a sequence 
``goes to'' a particular value. In this case we say the {\it limit\/} 
of the sequence exists.
Our goal will be to give a precise definition
of convergence. 

\example{This example gets 
close and then ``bounces away'', 
but eventually it ``goes to'' zero. We would say that this sequence 
converges to $0$.}
$$
1/1, \; 100, \;1/2, \; 100/2, \; 1/3, \; 100/3, \; 1/4, \; 100/4, 
\; \ldots
$$

\example{This example oscillates about $0$, but 
gets smaller and smaller in absolute
value, moving closer and closer to zero. In this case we would say 
the sequence converges to $0$.}
$$
-1, \; 1/2, \; -1/3, \; 1/4, \; -1/5, \; 1/6, \; -1/7, \; 1/8,
\; \ldots
$$

\example{This sequence has elements that get close
but then ``bounce away''.
Therefore, this sequence does not \index{converge} to $0$ as there is always an element
in the tail of the sequence which is ``far'' away from $0$. Nor can it
converge to $1$ (or to anything else) -- every tail also contains
elements $1/n$ which are far from $1$.}
$$
1, \; 1, \; 1/2, \; 1, \; 1/3, \; 1, \; 1/4, \; 1, \; 1/5, 
\; 1, \; 1/6, \; 1, \; \ldots
$$

\example{This sequence has elements that get close 
but then ``bounce away''.
Therefore, this sequence does not converge to zero as there is always 
an element in the tail of the sequence which is ``far'' away from 0.}
$$
1, \; 0.0001, \; 1/2, \; 0.0001, \; 1/3, \; 0.0001, \; 1/4, 
\; 0.0001, \; 1/5, \; 0.0001, \; \ldots
$$

\example{This sequence has elements that tend to $1$.
The farther out one goes in the sequence, the tail gets close and closer to
$1$. In this case we would say this sequence converges to $1$.}
$$
1 + 1, \; 1+ 1/2, 1 + 1/3, \; 1 + 1/4, 1 + 1/5, \; 1 + 1/6 \; \ldots
$$

In the examples where the sequence ``goes to''
a point we see that the tails 
get closer and closer to this point.
In the examples where there wasn't \index{convergence}
we saw that getting close is a relative statement. In example 5,
individual elements get arbitrarily close to $0$, but every tail still
contains elements a full $0.0001$ away -- the tail {\it as a whole\/}
never gets closer than $0.0001$ to $0$. So, it is not enough for {\it elements\/} of a sequence to
occasionally get closer to a limit point.
For convergence we need the {\it tails\/} to get closer to the limit point; and, not just closer, I think you will 
agree that we need the tails to
get {\it arbitrarily close\/} to the limit point. 
How do we precisely define the phrase {\it arbitrarily close\/}?

One way to say this is that a sequence converges to a number $a$ if 
the sequence eventually gets within {\it every\/} 
prescribed ``closeness'' tolerance of $a$.

This leads to a definition of \index{convergence} for infinite sequences.

\proclaim Definition. We say the sequence 
$\{x_n\}_{n=1}^{\infty}$ has the limit $a$, written
$\lim_{n \rightarrow \infty} x_n = a$, if for {\it every}
$\varepsilon > 0$ (closeness tolerance), there is a point
in the sequence, $N$, (which may depend on $\varepsilon$) at which all the
elements of its tail, $T_N$, are within $\varepsilon$ of $a$.

We now have a precise definition of what it means for a sequence to have a
limit, but how do we actually show a given sequence converges to a given limit
point? Notice work has to be done for {\it every} $\varepsilon > 0$.  The problem
is that there are an infinite number of such $\varepsilon$'s!

The only way to be able to handle an infinite collection of such $\varepsilon$'s is
to automate the process of finding an $N$ for a given $\varepsilon$. That is, 
we need to be able to find a function for $N$ in terms of $\varepsilon$. Then, 
given an $\varepsilon$, we use this function to produce an $N$ and then argue 
that any element in the corresponding tail is within this $\varepsilon$ of 
the limit value. In effect, we come up with one (or maybe a handful) of generic proofs
which cover all possible tolerances.

{\bf Note:\/} This is, in general, a quantum leap in difficulty from what 
you are probably used to 
seeing. Although I hope you have a good idea of the what and why of a limit, 
don't be discouraged if you find the next example and 
the theorem after that hard. It took mankind a long time to figure out this 
definition and use it. You should also realize this definition 
is based on the examples we considered. The definition is a formalization 
of an agreement that was felt to embody \index{convergence}; that is, it is 
a human activity. Definitions are man made constructions based on 
concrete problems.

\proclaim Worked Example. Prove that $\lim_{n \rightarrow \infty} x_n = 0$
when $x_n = 1/n$.

{\bf proof:\/} We need to write a program to handle {\it any\/} $\varepsilon$ 
as input. That is, if
someone were to hand us a \index{closeness tolerance}, 
$\varepsilon$, we would need to produce an $N$ such that all elements in the 
tail are within this $\varepsilon$ of $0$. This means all elements of the 
tail satisfy $|x_n - 0| < \varepsilon$, which is
$|x_n| < \varepsilon$; or, $x_n < \varepsilon$, since all
of the $x_n$ are positive.
What do we do? 

Let's first start with a particular $\varepsilon$ and see if we can handle that.
Then maybe we can see how to {\it automate\/} our procedure to handle 
{\it any\/} $\varepsilon$. If $\varepsilon$ is $1/10$, for instance, what would 
our choice be for $N$? Well, there are many choices for $N$. We are asking for a 
place to cut the sequence so that all elements of its tail are within $1/10$ 
of the limit point $0$.
We claim $N = 10$ works. Why? Because the tail in this case is:
$T_{10} = \{1/11, 1/12, 1/13, 
\ldots\,\}\!\!$~. It's clear that choosing $N$ to be $100$ or $137$ or 
$1000$ works as well.
This is because each of their tails: $T_{100} = \{1/101, 1/102, 1/103, \ldots\,\}$ or
$T_{137} = \{1/138, 1/139, 1/140, \ldots\,\}$ or
$T_{1000} = \{1/1001, 1/1002, 1/1003, \ldots\,\}$ all have the property that 
their elements are within $1/10$ of the limit point $0$. 
We only need to find {\it an\/} $N$ that works, 
let's pick $N$ to be $10$. So, we have found an $N$ for a {\it particular\/} 
$\varepsilon$. Can we automate this process so that we can
find an $N$ for {\it any\/} $\varepsilon$? Notice that we found our $N$ by
finding the first $N$ such that $1/(N+1)$ was less than $1/10$. If we do
something like this for an arbitrary $\varepsilon$ we will have handled the automation step.
A convenient choice for $N$ is\footnote{\kern 0.5pt \raise 0.5ex \hbox{\dag}}
{The floor function takes a number, $x$, and returns the greatest integer
which is less than or equal to $x$. In effect, it ``rounds down''
to the nearest integer.}
$N = {\rm floor}(1 / \varepsilon) + 1$. This is not always the {\it first\/}
$N$ that works -- for $\varepsilon = 1/10$ it gives $N = 11$ rather than
$10$ -- but it is a slightly generous choice which is always safe, and we
only need {\it an\/} $N$ that works. Try a few choices of $\varepsilon$ and see
if this makes sense. For this choice of $N$ every element in the tail $T_N$
is within $1/N$ of $0$ which is within $\varepsilon$ of $0$.

We have done what 
was asked: \goodbreak ``Given any 
closeness tolerance $\varepsilon>0$ find an $N$ (which usually depends on $\varepsilon$) 
such that all elements in the tail of 
the sequence are within $\varepsilon$ of the limit value.''

Proofs of this kind are probably different from what you are used to.
You are asked to automate a proof of a generic example; i.e., that of a given
$\varepsilon$. Further, the choice of $N$ is not unique. In the example above 
letting $N =100 * ({\rm floor} (1 / \varepsilon) + 1)$
is also a 
valid choice for $N$.
In addition, the way you
find this $N$ (based on a given $\varepsilon$) will change from proof to proof. 

One thing I think we all assume from our definition is that if a sequence
converges to a number $a$, then $a$ is unique. That is, a sequence can't
converge to two different numbers. This turns out to be true under our 
definition, but it requires a proof.

\proclaim Theorem(Limits are Unique). If the limit of a sequence exists, 
it is unique.

{\bf proof:\/} We will show that if the sequence $x_n$ converges to $a$, 
$\lim_{n \rightarrow \infty} x_n = a$, and $x_n$ converges to $b$, 
$\lim_{n \rightarrow \infty} x_n = b$, then $a$ and $b$ are equal.

We first sketch the outline of our argument. The key idea is that as we go 
farther and farther out in the sequence, the tail should be
very close to both $a$ {\it and\/} $b$; in fact {\it arbitrarily close\/}. 
This suggests that, in turn, 
$a$ and $b$ must be
very close; in fact {\it arbitrarily close\/}. But to be arbitrarily close 
means that 
they are the same number. Therefore, our strategy is to show that
$|a - b|$ is less than any prescribed tolerance. This will then mean
that $|a - b| = 0$; that is, they are the same number.
We now show that when $|a - b| < \varepsilon$ for {\it every\/}
$\varepsilon > 0$, then $a = b$. This is really the same as showing that when
a number $c$ satisfies $|c| < \varepsilon$ for {\it every\/}
$\varepsilon > 0$, then $c = 0$. 

{\bf claim:\/}
Suppose we have a non-negative number, $c$, which has the property that $c <
\varepsilon$ for all $\varepsilon > 0$, then $c = 0$. One might say that $c$ is 
arbitrarily close to $0$. If 
$c$ isn't 0, then let $\varepsilon = c/2$. In this case $c$ must be less than this $\varepsilon$, so
$c < c/2$. Since $c$ is positive the inequality is preserved after dividing by $c$ yielding
$1 < 1/2$. This is a contradiction; therefore $c$ must be zero. 

We now continue with our proof and show that $|a-b| < \varepsilon$ for all
$\varepsilon > 0$. Thus showing $a$ and $b$ are equal.

Let $\varepsilon>0$ be a given number, then since we know $x_n$ converges 
to $a$, given an $\varepsilon_1$ we know that there exists an $N_1$ such that 
$|x_n - a| < \varepsilon_1$ for all $x_n$ in
the tail $T_{N_1}$. Similarly, since $x_n$ converges to $b$, given
an $\varepsilon_2$
we know there exists an $N_2$
such that  $|x_n - b| < \varepsilon_2$ for all $x_n$ in
the tail $T_{N_2}$. If we let $N = \max(N_1,N_2)$, then the element
$x_{N+1}$ lies in both tails $T_{N_1}$ and $T_{N_2}$, so
$$
|a - b| \,= \,|a - x_{N+1} + x_{N+1} - b| \, \le \, |a - x_{N+1}| +
|x_{N+1} - b|  \, < \,
\varepsilon_1 + \varepsilon_2
$$

Now we need to choose $\varepsilon_1$ and $\varepsilon_2$ so that their sum is 
less than or equal $\varepsilon$. This will certainly be the case if both are 
chosen to be equal to $\varepsilon / 2$. 
So, for an arbitrary $\varepsilon>0$ the relation,
$|a-b| < \varepsilon$ holds. Therefore, $a$ is equal to $b$ 
and we have shown that 
the limit of a sequence is unique.

\section{Computational Strategy for Limits}
One of the things we would like to be able to do, just as was done when computing
derivatives, is to come up with ways to find the limits of sequences formed 
from simpler sequences. When computing derivatives we had a basic set of
functions and computed their derivatives from the definition. We then had 
rules for how to find 
the derivatives of functions that were composed of these basic functions. 
We would like to invoke a similar strategy when computing limits. To this 
end, we 
would like to compute the limit of a sequence by recognizing it as a ``base'' sequence or by
recognizing it as a composition of base sequences and use some set of rules to 
determine the limit in terms of limits of its base component sequences.
Hopefully, we could avoid using the definition directly to work 
out a given limit just as we were able to avoid using the direct definition of 
the derivative to compute derivatives of certain classes of functions.

In what follows we assume that $\{x_n\}_{n=1}^{\infty}$ and
$\{y_n\}_{n=1}^{\infty}$ are sequences with limits $a$ and $b$ 
respectively. 

Here are some rules for combining sequences stated without proof:

\proclaim Theorem(\index{Linearity of Limits}). If $z_n = c_1 x_n + c_2 y_n$ then 
$$
\lim_{n \rightarrow \infty} z_n = c_1 a + c_2 b
$$

\proclaim Theorem(\index{Product of Limits}). If $z_n = x_n * y_n$ then 
$$
\lim_{n \rightarrow \infty} z_n = a * b
$$

\proclaim Theorem(\index{Quotient of Limits}). If $z_n = x_n / y_n$ and 
$b \neq 0$ and $y_n \neq 0$ (for all $n$) then
$$
\lim_{n \rightarrow \infty} z_n = a / b 
$$

\proclaim Theorem(\index{Squeeze Limit}). If $x_n \le z_n \le y_n$ and $a = b$ then
$$
\lim_{n \rightarrow \infty} z_n = a = b
$$


To this we add a specific set of sequences for which we can
compute the limit. We will need the notion of a {\it bounded\/} sequence:
a sequence is bounded if there is a number $M$ such that
the absolute value of every element in the sequence is less than
or equal to $M$.


\proclaim Theorem(\index{Constant Limit}). If $x_n = k$ for some number $k$, then 
$$
\lim_{n \rightarrow \infty} x_n = k
$$

\proclaim Theorem(\index{Generalized Decreasing Sequence}). If $x_n = b_n / n$, where $\{b_n\}_{n=1}^{\infty}$ is a bounded sequence, then
$$
\lim_{n \rightarrow \infty} x_n = 0
$$

\proclaim Theorem(\index{Power Limit}). If $x_n = k^n$ with $|k| < 1$, then
$$
\lim_{n \rightarrow \infty} x_n = 0
$$


As an example of the limit of the generalized decreasing \index{sequence} pattern consider
the sequence 
$\{x_n\}_{n=1}^{\infty}$,
where $x_n = \sin(n) / n$. In this case $b_n = \sin(n)$. The only thing we 
need to do is check that $|b_n|$ is bounded. This is easy as the 
absolute value of the $\sin$ of anything is less than or equal to  $1$, 
for instance. Therefore,
$$
\lim_{n \rightarrow \infty} \sin(n)/ n = 0
$$

As another example, consider the sequence $z_n = 1/n + \sin(n)/n$.
We can view $z_n$ as a sum of two sequences $1/n$ and 
$\sin(n)/n$. Since each of these sequences have limits 
their sum does also and the value of the limit is the sum of 
the limits of $1/n$ and $\sin(n)/n$. Therefore, 
$$
\lim_{n \rightarrow \infty} z_n = \lim_{n \rightarrow \infty} 1/n + 
\lim_{n \rightarrow \infty} \sin(n)/n = 0 + 0 = 0
$$
Finally, if a given sequence, $\{x_n\}_{n=1}^\infty$, converges to $a$, is it
the case that $\{f(x_n)\}_{n=1}^\infty$ converges to $f(a)$? This is a nice computational property to
have as its says: $\lim_{n\rightarrow \infty} f(x_n) = f(\lim_{n \rightarrow \infty} x_n) = f(a)$. In other words, 
if a function has this property, limits ``pass through'' it.
This property holds for 
``nice'' functions. In fact, we will call a function {\it continuous\/} at $a$ 
if this property holds. Functions that are continuous at every input value in their domain are called continuous functions.

So, what do \index{continuous function}s look like and how 
can one spot them. 
They have the property that while drawing their graph, one would not have to 
lift the pen off the page.
All polynomials are continuous for all inputs; the logarithm functions 
are continuous for positive inputs; the \index{exponential function}s, $a^x$, for $a>0$ 
are all continuous for all inputs. Sums, products, and differences of continuous 
functions are continuous. Ratios of continuous functions are continuous provided the denominator 
stays away from 0. For an example of a discontinuous function see
example 3 in appendix A.

One more fact ties continuity to the derivative: a function which is
{\it differentiable\/} at $a$ is automatically {\it continuous\/} at $a$.
Why? Let $\{x_n\}_{n=1}^\infty$ be a sequence converging to $a$. For the
terms with $x_n \ne a$ we may write
$$
f(x_n) - f(a) = \biggl( {f(x_n) - f(a) \over x_n - a} \biggr) (x_n - a)
$$
The first factor is an approximating slope, which tends to $f'(a)$ since $f$
is differentiable at $a$; the second factor tends to $0$. By the Product of
Limits Theorem their product tends to $f'(a) * 0 = 0$. (For any terms with
$x_n = a$ the difference $f(x_n) - f(a)$ is already exactly $0$.) That is,
$\lim_{n \rightarrow \infty} f(x_n) = f(a)$ -- which is precisely the
definition of continuity at $a$.

\section{Convergence Based on Properties of a Sequence}
We have seen how to show a sequence converges to a given value. What if 
we don't have a candidate for the limit? Can we tell if a sequence will 
converge just by examining properties of the sequence itself? The answer is 
yes, in certain situations we know that a limit of the sequence exists.

\proclaim Theorem (\index{Monotone Sequence}). Suppose $\{x_n\}_{n=0}^\infty$ is a 
sequence which is bounded above by $M$; meaning, $x_n \le M$ for all $n$.
If the sequence is also non-decreasing ($x_{n+1} \ge x_n$ for all $n$) then
the limit of the sequence exists. Similarly, if a sequence is bounded 
below and is non-increasing it also has a limit.

There are other criteria, but they are outside the scope of this book.

\section{Two Trigonometric Limits}
In the chapter on Differential Calculus, computing the derivative of the
sine function came down to two limits: for any sequence
$\{h_n\}_{n=1}^\infty$ which converges to $0$ with the property that
$h_n \ne 0$ for all $n$,
$$
\lim_{n\rightarrow \infty} \frac{\sin(h_n)}{h_n} = 1
\qquad {\rm and} \qquad
\lim_{n\rightarrow \infty} \frac{\cos(h_n) - 1}{h_n} = 0
$$
In the chapter these were taken on faith; here we pay the debt. The first
limit is a fine application of the \index{Squeeze Limit} Theorem of this
appendix; the second follows from the first.

Examining the figure below and comparing the areas of the triangular region ABC to 
the sector ADC to the triangle ADE we find that 
$$
\obrace{\frac{\cos(\theta)\sin(\theta)}{2}}{Area of triangle ABC} \le 
\obrace{\frac{\theta}{2}}{Area of sector ADC} \le 
\obrace{\frac{1 * \tan(\theta)}{2}}{Area of triangle ADE}
$$
\epsfigure{eps/sinderiv.eps}{Comparison of triangle ABC, sector ADC, 
and triangle ADE.}
Here, we are assuming $\theta$ is positive, measured in radians, and 
is less than $\pi/2$. The area of each triangle is found using the formula, 
base * height / 2, while the area of the sector is determined by the formula 
$\theta r^2/2$. In our case, $r = 1$, giving $\theta/2$; again, 
provided $\theta$ is measured in \index{radian}s. We briefly describe 
radian measure for those who have not seen it before.

The radian measure of an angle is the arc 
length of the angular sector 
formed by this angle on the unit circle. That is, if a sector has angle $\theta$ in 
radians, then the corresponding arc (on the unit circle) has length $\theta$ (see 
the figure above).
The arc length of a complete circle 
(360 degrees) is $2\pi$; therefore, there are $2\pi$ radians in 360 degrees. This means
that there are ${360 \over 2 \pi }= {180 \over \pi}$ (which  is approximately 57) degrees per radian. 
Also, the area of the unit circle is $\pi$; so, 
the area of a sector of $\theta$ radians (on the unit circle) is a fraction of the area of the 
full circle $\pi$. What fraction? Well, the sector of $\theta$ radians has an arc length of $\theta$ while
the arc length of a complete circle is $2 \pi$. Since the area of a complete circle is $\pi$, the area
of the sector must be ${\theta \over 2 \pi}$ of $\pi$. So, the area of the sector is $\theta / 2$.
One can use other angular measures (such as degrees or grads) 
but the formulas for the area of a sector have messy constants; and, in turn, 
the derivatives of the trig functions using these angular measures also 
have messy constants. Note that this is not unlike the derivative 
of the logarithmic functions; any choice of base other than $e$ leads to a messier 
formula.

Continuing with our calculations, the first inequality gives 
the relation:
$$
\frac{\sin(\theta)}{\theta} \le \frac{1}{\cos(\theta)}
$$
while the second inequality yields the relation:
$$
\cos(\theta) \le \frac{\sin(\theta)}{\theta}
$$
Combining these gives:
$$
\cos(\theta) \le \frac{\sin(\theta)}{\theta} \le \frac{1}{\cos(\theta)}
$$
A similar argument shows that the relation holds when $\theta$ is negative and larger 
than $-\pi/2$. We may replace $\theta$ with the $n^{\rm th}$ element of our sequence $h$ giving:
$$
\cos(h_n) \le \frac{\sin(h_n)}{h_n} \le \frac{1}{\cos(h_n)}
$$
Now, since the limit of the first and third terms exists and is $1$, it is the case that the middle term has 
the same limit.\footnote{\kern 0.5pt \raise 0.5ex \hbox{\dag}}{This follows from the Squeeze Limit Theorem of this appendix.} Therefore,
$$
\lim_{n\rightarrow \infty} \frac{\sin(h_n)}{h_n} = 1
$$
We can use this result to compute the other limit of interest, 
$(\cos(h_n) - 1)/h_n$. First note that 
$$
\frac{\cos(h_n) - 1}{h_n} = \frac{(\cos(h_n) - 1) (\cos(h_n) + 1)}{h_n (\cos(h_n) + 1)}
= \frac{\cos^2(h_n) - 1}{h_n(\cos(h_n) + 1)}
$$
This may be written 
$$
\frac{-\sin^2(h_n)}{h_n(\cos(h_n) + 1)} = -\frac{\sin^2(h_n)}{h_n^2} \frac{h_n}{\cos(h_n) + 1}
$$
Since the last expression is a product of expressions which have limits, we may 
write:\footnote{\kern 0.5pt \raise 0.5ex \hbox{\ddag}}{This follows from the Product of Limits Theorem of this appendix.}
$$
\lim_{n\rightarrow \infty}\frac{\cos(h_n) - 1}{h_n} = -\lim_{n\rightarrow \infty}
\frac{\sin^2(h_n)}{h_n^2} \lim_{n\rightarrow \infty}\frac{h_n}{\cos(h_n) + 1}
= -(1^2)
 * (0 / 2) = -1 * 0 = 0
$$

These are the two limits used in the chapter on Differential Calculus to
show that $\sin'(x) = \cos(x)$ and $\cos'(x) = -\sin(x)$.

\section{Sequences without Limits}
Not all sequences have limits, in fact most do not. 
In this section we give examples of sequences which don't have limits.

\example{Let $x_n = (-1)^n$. This sequence bounces between $1$
and $-1$ and so doesn't converge to anything.}

\example{Let $x_n = n$. This sequence gets larger and larger and 
so it cannot converge to anything.}

\example{Let $x_n = \sin(n)$. This sequence bounces around the interval $[-1,1]$ 
but doesn't converge to anything.}


